\documentclass[11pt]{article}
\usepackage[spanish,mexico]{babel}
\usepackage[utf8]{inputenc}
\usepackage[T1]{fontenc}
\usepackage{geometry}
\usepackage{hyperref}
\usepackage{enumitem}
\geometry{margin=2.5cm}
\title{Estado del arte: control de misiones agenticas, trazabilidad y supervision humana}
\author{Sistema Operativo de Tesis de Posgrado}
\date{31 de mayo de 2026}
\begin{document}
\maketitle

\section*{Resumen}
Los sistemas multiagente recientes desplazan el problema desde ``generar una respuesta'' hacia coordinar tareas, herramientas, memoria, permisos y evidencia. En un sistema de tesis, el control de misiones debe priorizar trazabilidad y control humano sobre autonomia total.

\section*{Hallazgos}
\begin{enumerate}[leftmargin=*]
\item La verificacion multiagente y los patrones de test-time compute muestran valor al separar generadores, revisores y verificadores.
\item Las arquitecturas agenticas modernas requieren limites explicitos: permisos por herramienta, logs, presupuestos, timeouts y recuperacion ante fallo.
\item La supervision humana debe quedar modelada como protocolo, no como comentario: una tarea puede ser auditada por IA, pero la validacion formal requiere consentimiento humano y Step ID.
\item Mission Control debe almacenar misiones, estados, evidencias, artefactos y bloqueos sin ocultar incertidumbre ni convertir resultados parciales en cierre.
\end{enumerate}

\section*{Implicaciones para el Centro de Control}
El Centro de Control debe operar como interfaz soberana: recibir misiones, crear entregables verificables, disparar pruebas, registrar evidencia y presentar bloqueos. Para trabajos academicos, cada mision debe terminar con rutas de \texttt{.tex}, \texttt{.pdf}, fuentes usadas, log de compilacion y estado de auditoria.

\section*{Referencias seleccionadas}
\begin{itemize}[leftmargin=*]
\item Multi-Agent Verification: Scaling Test-Time Compute with Multiple Verifiers, arXiv:2502.20379.
\item MetaGPT: Meta Programming for a Multi-Agent Collaborative Framework, arXiv.
\item Stanford CRFM HELM, enfoque de evaluacion transparente y reproducible.
\item ISO/IEC 25059:2023, modelo de calidad para sistemas de IA.
\end{itemize}
\end{document}
